% ---------------------------------------------------------------------------
% Paper-format draft for supervisor review. Single file; compiles on Overleaf
% with pdflatex (run twice for references).
%
% Every number in this draft is taken from committed logs on branch
% fix/post-close-carry-forward (see runs/README.md and docs/RESEARCH_LOG.md).
% Figures marked "exploratory" were not pre-registered.
%
% LaTeX notes carried over from data_window_memo.tex:
%   1. \textcent is text-mode only; use \cent (defined below) in maths.
%   2. Maths inside \textbf{} is not bolded; use \bm{}.
%   3. Check WARNINGS as well as errors when building.
% ---------------------------------------------------------------------------
\documentclass[11pt,a4paper]{article}

\usepackage[T1]{fontenc}
\usepackage[utf8]{inputenc}
\usepackage{textcomp}
\usepackage{amsmath}
\usepackage{amssymb}
\usepackage{bm}
\usepackage{booktabs}
\usepackage{multirow}
\usepackage[margin=2.4cm]{geometry}
\usepackage{enumitem}
\usepackage{xcolor}
\usepackage{url}
\usepackage{hyperref}
\hypersetup{colorlinks=true, linkcolor=black, citecolor=blue!50!black, urlcolor=blue!60!black}

\newcommand{\cent}{\text{\textcent}}
\newcommand{\x}{\ensuremath{\times}}
\setlist[itemize]{itemsep=2pt, topsep=4pt}
\setlist[enumerate]{itemsep=2pt, topsep=4pt}

\title{\textbf{When Does Graph Learning Beat a No-Arbitrage Constraint?}\\[4pt]
\large Coherence and Learnability in Kalshi Event Contracts}
\author{[Author name]\\[2pt]
\small Final Year Project --- draft for supervisor review}
\date{\today}

\begin{document}
\maketitle

\begin{abstract}
Mutually exclusive, collectively exhaustive (MECE) baskets of event contracts must
be priced to sum to \$1, and threshold ladders must be monotone in strike. These
no-arbitrage constraints give free structure that a learned model could exploit.
We ask where a spatio-temporal graph attention network (STGAT) adds value over the
constraint itself, and where the constraint suffices. Using Kalshi trade data and
hourly top-of-book quotes, we find three things.
(i) \emph{The baskets are priced coherently.} Where every leg has a two-sided
quote, the mid-price sum is within $0.34\cent$ of \$1. The apparent dislocation
in last-trade prices is almost entirely microstructure: 82--87\% comes from the
side of the book on which trades print, and the rest from one-sided books.
(ii) \emph{The constraint is a strong forecaster, and how strong depends on its
form.} Shrinking toward \$1 cuts basket forecast error by a third against
persistence (shrinkage factor $0.55$). The ladder's inequality constraint gives
almost no reversion ($0.94$).
(iii) \emph{The STGAT learns traded-price dynamics well, but the skill stops at
executable prices.} It forecasts basket trade prices at $0.61\x$ persistence
error, and learns settlement convergence without being given time to close; the
graph is what lets it compete with the constraint. On the executable bid,
however, it is $1.82\x$ the constraint oracle. Its entry timing beats a coin flip
but not immediate entry, and no strategy is profitable after fees at quoted
prices. The learned model adds value in predicting how prices \emph{trade}; the
constraint suffices for what they are \emph{worth}. We also document four
evaluation pitfalls that produced false results during the project, and a test
that detects each.
\end{abstract}

% ===========================================================================
\section{Introduction}
% ===========================================================================

Prediction markets list contracts whose payoffs are tied together by construction.
In a MECE basket, such as ``Which range will New York's high temperature fall in
tomorrow?'', exactly one leg pays \$1, so the legs' prices should sum to \$1. In a
threshold ladder (``above $K_1$'', ``above $K_2$'' with $K_1<K_2$), the higher
strike cannot be worth more than the lower one. These no-arbitrage constraints are
known exactly, before any data is seen.

This makes event contracts an unusual test bed for learned forecasting models. In
most financial settings, a model must discover cross-asset structure from noisy
data. Here the structure is given, and the question becomes sharper:
\begin{quote}
\emph{Where does a learned model add value over the constraint itself, and where
does the constraint suffice?}
\end{quote}

We study this question on Kalshi, a CFTC-regulated exchange, using a
spatio-temporal graph attention network (STGAT) whose graph encodes the
constraints: basket legs are linked through a basket hub, and ladder strikes are
linked in order. We compare it against baselines that use the constraint alone,
at two kinds of price: the last traded price, which most studies use, and the
executable top-of-book quote, which is what a trader would actually pay or
receive.

\paragraph{Contributions.}
\begin{enumerate}
  \item \textbf{A coherence measurement.} Using hourly best bid/ask quotes for
    10{,}981 basket legs, we show that Kalshi baskets are priced coherently at the
    mid. We decompose the apparent last-price dislocation into bid/ask print side
    and one-sided books. The size of the print-side effect was predicted before it
    was measured (Section~\ref{sec:coherence}).
  \item \textbf{The constraint as a forecaster.} A single shrinkage factor toward
    \$1 is a strong baseline for baskets and a weak one for ladders. We explain the
    difference by the constraint's form: an equality pins a level, while an
    inequality only bounds a half-line (Section~\ref{sec:constraint}).
  \item \textbf{Where learning helps.} An ablation ladder runs from persistence,
    through constraint-only and family-aware shrinkage and a graph-free temporal
    model, to the full STGAT. It locates the model's value in traded-price
    dynamics and settlement convergence, and shows that the graph is necessary
    for it (Section~\ref{sec:learning}).
  \item \textbf{Executable boundaries.} At quoted prices, the model has no
    forecasting skill on the bid, and its timing beats chance but not immediate
    entry. No strategy is profitable after fees (Section~\ref{sec:exec}).
  \item \textbf{Evaluation pitfalls.} Four traps, each of which produced a false
    result during this project, together with the test that detects each
    (Section~\ref{sec:pitfalls}).
\end{enumerate}

% ===========================================================================
\section{Background}
% ===========================================================================

\paragraph{Prediction markets.} Prices of binary event contracts are commonly read
as probabilities \cite{wolfers2004}. They are generally informative, but show
systematic biases, most prominently the favourite--longshot bias
\cite{snowberg2010}. On Kalshi specifically, contract prices become more accurate
as markets approach close. Takers lose about 32\% on average and makers about
10\% \cite{burgi2025}.

\paragraph{Basket dislocations.} Saguillo et al.~\cite{saguillo2025} estimate about \$40M of realised profit from MECE baskets on
Polymarket priced away from \$1 (April 2024--April 2025). Polymarket was largely
fee-free for trading over that period, and the profits went to automated traders
acting on the live order book. Whether the same dislocations are executable on a
fee-charging, market-maker-intermediated venue is an open empirical question,
and one we address.

\paragraph{Microstructure of trade prices.} A trade executes at either the bid or
the ask, so trade prices ``bounce'' across the spread even when the efficient price
does not move \cite{roll1984}. Trade and quote series can be decomposed into an
efficient price plus microstructure noise \cite{hansen2006}. Trade direction is
classically inferred from quotes \cite{lee1991}, with about 85\% accuracy
\cite{oddswhite2000}. Kalshi instead reports the aggressor side of every trade
directly. Constraints spanning several instruments are further affected by
non-synchronous trading \cite{lomackinlay1990}. Classic put--call parity tests
restrict comparisons to near-simultaneous prices for this reason
\cite{klemkosky1979, kamara1995}.

\paragraph{Graph neural networks for markets.} Graph attention
\cite{velickovic2018} and spatio-temporal graph models \cite{yu2018stgcn} combine
cross-sectional structure with temporal dynamics. In market microstructure,
deep models usually forecast \emph{changes} in a reference price, not levels
\cite{zhang2019deeplob}. We use a graph whose edges are the known constraints,
not a learned adjacency.

% ===========================================================================
\section{Motivation}
% ===========================================================================

Early in this project, a model trained on last-trade prices appeared to beat an
optimal constraint-based shrinkage forecast by 15\%. It also appeared to identify
tradeable dislocations. Both claims failed under scrutiny, for reasons that
motivate the design of this paper.

\begin{itemize}
  \item \textbf{Last-trade basket sums are dominated by microstructure.} 68.7\% of
    basket trade prints are ask-side. With a half-spread of $2.8\cent$, that
    predicts a $+6.3\cent$ mean excess on a 6-leg basket. The measured excess was
    $+5.25\cent$. A model that forecasts last-trade sums may therefore be
    forecasting which side the next prints land on, not mispricing.
  \item \textbf{Several pipeline leaks inflated the result.}
    \begin{itemize}
      \item resolved markets were replayed after close;
      \item time-to-close used after-the-fact close times;
      \item basket membership was selected with hindsight;
      \item baskets were scored on a subset of their legs.
    \end{itemize}
    Section~\ref{sec:pitfalls} quantifies each.
\end{itemize}

We therefore (a) rebuilt the pipeline causally and pre-registered every re-run
before seeing results; (b) scored forecasts against baselines that use the
constraint, not only against persistence; and (c) measured executable prices
directly from quotes rather than inferring them from trades.

% ===========================================================================
\section{Method}
% ===========================================================================

\subsection{Population and graph}
Contracts are nodes on a 2-hour snapshot grid. A \emph{MECE basket} is an event
whose legs partition the outcome space. Membership is decided causally in
validation and test: a family qualifies only on instances known to have resolved
at least 24h earlier, and the event's own outcome is never used. (Training months
use a curated membership list.) Each basket has
a hub node, linked to its legs by directed edges in both directions. \emph{Ladder}
pairs are adjacent strikes of the same underlying, linked by a monotonicity edge.
Node features summarise each 2-hour window: last price, VWAP, return, volatility,
volume, order flow and hub leg counts. Time-to-close is zeroed as an input in
every split, and ladder edge features are zeroed.

\subsection{Model}
The STGAT has four parts:
\begin{itemize}
  \item a type-specific input projection;
  \item a two-channel spatial attention layer, whose channels are MECE hub--leg
    and ladder;
  \item per-node \emph{causal} temporal self-attention over the snapshot sequence;
  \item a residual forecast head, $\hat p_{t+h} = p_t + \Delta_h$, with the final
    layer zero-initialised so training starts exactly at persistence.
\end{itemize}
Horizons are $h\in\{1,2,3,6,12\}$ snapshots, i.e.\ 2--24 hours. The loss combines
per-node error with a basket-sum term; the ladder term's weight is zero.
Checkpoints are selected on validation MECE error.

\subsection{Targets and baselines}
For basket $b$ with legs $L_b$, let $S_t=\sum_{i\in L_b} p_{i,t}$. Under the
\texttt{whole\_settle} rule, a basket is scored only if every listed leg is
present. A leg that settles by $t+h$ is valued at its settlement (0 or 1).
Baselines:
\begin{itemize}
  \item \textbf{Persistence:} $\hat S_{t+h}=S_t$.
  \item \textbf{Constraint oracle:} $\hat S_{t+h}=1+\alpha(S_t-1)$, with the single
    $\alpha$ that minimises squared error \emph{on the scored split}. This is
    deliberately optimistic.
  \item \textbf{Family oracle:} the same, with one $\alpha$ per series \x\
    horizon. A basket's family is known in advance.
\end{itemize}

\subsection{Executable prices}
Hourly best bid and ask for every basket leg come from Kalshi's public historical
candlestick endpoint \cite{kalshiapi}, aligned to the snapshot window ends.
A bid of 0, or an ask of \$1, is a real one-sided book and is kept as the
executable price.
\begin{itemize}
  \item \textbf{Bid-side forecast.} The model's predicted change in the
    last-trade sum is applied to the quoted bid sum:
    $\hat B_{t+h} = B_t + (\hat S_{t+h}-S_t)$. The half-spread cancels in the
    change if the spread is stable \cite{roll1984}.
  \item \textbf{Hold to settlement.} Selling a basket at the bids yields gross
    $(\sum_i \mathrm{bid}_i - 1)$ per set, and buying at the asks yields
    $(1-\sum_i \mathrm{ask}_i)$. Kalshi's taker fee is charged on the executed
    price.
\end{itemize}

\subsection{Timing control}
The model chooses to enter now, or to wait one step if it predicts the
dislocation will widen. The control makes the same choice by coin flip, on
identical events and sizes, so any gap is due only to the forecast. Confidence
intervals come from a cluster bootstrap over baskets.

% ===========================================================================
\section{Experimental Setup}
% ===========================================================================

\begin{itemize}
  \item \textbf{Data.} Kalshi trades and markets, May--September 2025, with a
    strictly forward split: train May--June, validation July--August, test
    September. Hourly quotes cover all 10{,}982 validation/test basket legs;
    10{,}981 were retrieved.
  \item \textbf{Seeds and reporting.} Four seeds per configuration. We report
    the mean and range over seeds, with a 95\% cluster-bootstrap CI per seed
    and the cluster count $k$.
  \item \textbf{Pre-registration.} Every re-run reported here was pre-registered
    in the repository before results were seen: the protocol, the decision rule
    and a stated prediction. The test split was scored once per seed per
    protocol. Deviations are disclosed.
  \item \textbf{Test positions.} 37{,}392 basket-horizon positions ($k=348$
    baskets), of which 18.6\% settle within the horizon.
\end{itemize}

% ===========================================================================
\section{Results}
% ===========================================================================

\subsection{Baskets are coherent at quoted prices}\label{sec:coherence}

\begin{table}[h]
\centering
\caption{Basket sums on the test split: whole baskets with every leg live.
Quotes cover 98.1\% of leg-snapshots and 97.1\% of basket-snapshots.}
\label{tab:coherence}
\begin{tabular}{lcc}
\toprule
 & all baskets ($n=7{,}532$) & every leg two-sided (69\%) \\
\midrule
Sum of last trades & $1.0405$ & $1.0261$ \\
Sum of quoted mids & $1.0185$ & $\bm{1.0034}$ \\
Sum of bids / asks & $0.944$ / $1.093$ & $0.967$ / $1.039$ \\
Sell at bids beats \$1 & 1.8\% of snapshots ($\approx 2\cent$) & --- \\
\bottomrule
\end{tabular}
\end{table}

Where every leg has a two-sided quote, the mid-price sum is within $0.34\cent$
of \$1 (Table~\ref{tab:coherence}). The same holds on validation ($+0.50\cent$).
In these baskets, 82--87\% of the last-trade excess is print side. Across all
baskets, the remaining excess sits in legs with no bid, whose mid is inflated.
Executable dislocation is rare and small. Selling at the bids beats \$1 in 1.8\%
of snapshots, by about $2\cent$, before a taker fee of $1.3$--$5.8\cent$ per
basket. The quoted per-leg spread (median $2.0\cent$, mean $2.9\cent$) is below
our earlier trade-based estimates ($4.7$--$5.6\cent$). The trade-print bid proxy
overstated quoted bids by $1.1\cent$ per leg on average.

\subsection{The constraint as a forecaster}\label{sec:constraint}

On test, the constraint oracle has 0.67\x\ the error of persistence for basket
sums ($\alpha=0.55$), but 0.985\x\ for ladder gaps ($\alpha=0.94$). Basket
deviations revert strongly to the level the constraint pins. Ladder gaps are only
bounded above by zero, so there is no level for them to revert to.

\subsection{Where learning helps: the ablation ladder}\label{sec:learning}

\begin{table}[h]
\centering
\caption{Test error, relative to the single-$\alpha$ constraint oracle unless
stated. Mean over 4 seeds (range). ``Unsettled'' = baskets with no leg settling
by $t+h$.}
\label{tab:ladder}
\begin{tabular}{lccc}
\toprule
Model & all positions & unsettled & vs.\ persistence \\
\midrule
Persistence & $1.49$ & --- & $1.00$ \\
Constraint oracle (single $\alpha$) & $1.00$ & $1.00$ & $0.67$ \\
Temporal model, no graph & $1.128$ & $1.106$ & $0.760$ \\
STGAT (full) & $\bm{0.909}$ ($0.883$--$0.941$) & $1.034$ & $\bm{0.613}$ \\
STGAT, unsettled-only retrain\textsuperscript{a} & --- & $0.950$ & $0.781$ \\
\midrule
\multicolumn{4}{l}{\emph{Against the family oracle (per series \x\ horizon), unsettled:}} \\
STGAT, unsettled-only retrain & --- & $1.052$ ($1.024$--$1.074$) & --- \\
Temporal model, no graph & --- & $1.224$ & --- \\
\bottomrule
\multicolumn{4}{l}{\footnotesize \textsuperscript{a} Settled baskets excluded from the loss and from selection.}
\end{tabular}
\end{table}

Three findings (Table~\ref{tab:ladder}):
\begin{itemize}
  \item \textbf{The STGAT forecasts traded basket prices well}, at $0.61\x$
    persistence error, consistently across seeds.
  \item \textbf{The graph is necessary.} Without the MECE and ladder edges, the
    model falls behind even the single-$\alpha$ constraint ($1.128\x$).
  \item \textbf{On test, the margin over the constraint comes from settlement
    convergence.} Baskets that settle by $t+h$ have a target of exactly \$1, and
    the model anticipates this even though time-to-close is not an input.
    Perturbation tests confirm that predictions at $t$ are unchanged when future
    snapshots or time-to-close are randomised, so this is not leakage. The
    model's pull toward \$1 rises smoothly as close approaches, with no jump at
    the settlement boundary, which points to family-level lifecycle learning.
    On unsettled baskets the model matches, but does not beat, the constraint
    plus contract family. On validation it does beat it (0.887\x\ single-$\alpha$
    unsettled), but that is the selection split.
\end{itemize}

\paragraph{Ladders (exploratory).} On 99{,}646 test ladder positions, the model
scores $0.982$--$0.993\x$ the oracle and $0.966$--$0.977\x$ persistence. The
graph-free model is marginally better ($0.974$--$0.982\x$), so the small ladder
gain comes from per-contract temporal forecasting, not from ladder structure.

\subsection{Executable boundaries}\label{sec:exec}

\begin{table}[h]
\centering
\caption{Executable-price results on test (4 seeds, cluster CIs by basket).}
\label{tab:exec}
\small
\begin{tabular}{p{7.4cm}cc}
\toprule
Quantity & result & CIs excluding null \\
\midrule
Bid-side forecast vs.\ oracle, unsettled ($n=30{,}001$, $k=344$) & $1.822\x$ ($1.80$--$1.85$) & 4/4 above 1 \\
Timing: model $-$ random, last prices (\$/trade, $k=312$) & $+0.192$ ($0.179$--$0.211$) & 4/4 above 0 \\
Timing: model $-$ random, quoted prices & $+0.867$ & 3/4 above 0 \\
Timing: model $-$ enter now, quoted prices & $-0.957$ & 4/4 \emph{below} 0 \\
Fee-aware entry at quotes: trades fired / net & 1 of 2{,}090 / $+\$0.07$ & --- \\
Static entry at quotes: net per trade & $-\$2.18$ & --- \\
\bottomrule
\end{tabular}
\end{table}

Three results (Table~\ref{tab:exec}):
\begin{itemize}
  \item \textbf{No bid-side skill.} Applied to the quoted bid, the model's
    forecast is $1.82\x$ the constraint oracle, and worse than persistence. Its
    predicted change tracks where trades print, not where the bid moves.
  \item \textbf{Timing beats chance, but not the trivial policy.} The model's
    entry timing beats a coin flip at last prices (4/4 seeds) and at quoted
    prices (3/4). At quoted prices, however, waiting is costly, and the model
    loses to entering immediately on every seed.
  \item \textbf{No profit after fees at executable prices.} A rule that enters
    only when the quoted edge exceeds the fee fired once in 2{,}090
    opportunities. This agrees with evidence that Kalshi takers lose on average
    \cite{burgi2025}. It contrasts with the fee-free Polymarket evidence
    \cite{saguillo2025}.
\end{itemize}

\subsection{Evaluation pitfalls}\label{sec:pitfalls}

\begin{table}[h]
\centering
\caption{Traps that produced false results during the project, and the test
that detects each.}
\label{tab:pitfalls}
\begin{tabular}{p{4.1cm}p{5.6cm}p{4.3cm}}
\toprule
Trap & Effect & Detection \\
\midrule
Settled contracts replayed after close & 57\% of old ladder test pairs were
frozen replays, producing ``ladders are unforecastable'' & audit of the share of
pairs past close \\
Settlement-pinned targets & all of the 0.909\x\ test margin came from baskets
settling within $h$ & stratify by settlement \\
Last trades treated as executable & bids overstated by $1.1\cent$ per leg;
spreads by $\approx$40\%; bid-side skill assumed & validate against quoted top
of book \\
Ablation scored with the wrong graph & graph-free checkpoints would have been
scored \emph{with} the graph & record dropped edges; known-answer test \\
\bottomrule
\end{tabular}
\end{table}

These were caught by reproduction guards, which check that every re-score
reproduces the committed figures exactly, and by pre-registration
(Table~\ref{tab:pitfalls}). One headline figure had no log at all, and was
recovered and found to have been computed on the selection split. The
provenance of every reported number is now logged.

% ===========================================================================
\section{Discussion and Limitations}
% ===========================================================================

\begin{itemize}
  \item \textbf{Answer to the question.} Graph learning adds value on traded-price
    dynamics: where trades print, and how baskets converge toward settlement. The
    graph is necessary for this. Wherever money changes hands, the constraint
    suffices. This follows from the coherence finding: if baskets are coherent at
    executable prices, no model should profit there.
  \item \textbf{Limitations.}
    \begin{itemize}
      \item Top-of-book quotes only, at hourly resolution; no depth, so fills at
        size are assumed.
      \item A single test month, scored for several model variants; the
        multiple-look caveat applies.
      \item The timing control's wait rate is not matched to the model's.
      \item Ladder figures are exploratory.
      \item The pilot window is five months, and the population is concentrated
        in weather and sports.
    \end{itemize}
  \item \textbf{Next.} A pre-registered 7-fold walk-forward over November 2024 to
    November 2025, with the protocol frozen from this study, fresh held-out
    months, and quotes for all months.
\end{itemize}

% ===========================================================================
\section{Conclusion}
% ===========================================================================
Kalshi's MECE baskets are priced coherently once microstructure is accounted for.
A no-arbitrage constraint is a strong forecaster where it pins a level, and a weak
one where it only bounds an inequality. A constraint-structured STGAT learns
traded-price dynamics well, and needs its graph to do so. That skill does not
extend to executable prices, where the constraint alone is as good or better.

% ===========================================================================
\begin{thebibliography}{99}
\bibitem{burgi2025} C.~B\"urgi, W.~Deng and K.~Whelan. Makers and Takers: The Economics of the Kalshi Prediction Market. CESifo / CEPR Discussion Paper DP20631, 2025.
\bibitem{hansen2006} P.~R.~Hansen and A.~Lunde. Realized variance and market microstructure noise. \emph{Journal of Business \& Economic Statistics}, 24(2):127--161, 2006.
\bibitem{kalshiapi} Kalshi. Get Market Candlesticks (historical). API documentation, \url{https://docs.kalshi.com/api-reference/market/get-market-candlesticks}, accessed 2026-09-27.
\bibitem{kamara1995} A.~Kamara and T.~W.~Miller. Daily and intradaily tests of European put-call parity. \emph{Journal of Financial and Quantitative Analysis}, 30(4):519--539, 1995.
\bibitem{klemkosky1979} R.~C.~Klemkosky and B.~G.~Resnick. Put-call parity and market efficiency. \emph{Journal of Finance}, 34(5):1141--1155, 1979.
\bibitem{lee1991} C.~M.~C.~Lee and M.~J.~Ready. Inferring trade direction from intraday data. \emph{Journal of Finance}, 46(2):733--746, 1991.
\bibitem{lomackinlay1990} A.~W.~Lo and A.~C.~MacKinlay. An econometric analysis of nonsynchronous trading. \emph{Journal of Econometrics}, 45:181--212, 1990.
\bibitem{oddswhite2000} E.~R.~Odders-White. On the occurrence and consequences of inaccurate trade classification. \emph{Journal of Financial Markets}, 3:259--286, 2000.
\bibitem{roll1984} R.~Roll. A simple implicit measure of the effective bid-ask spread in an efficient market. \emph{Journal of Finance}, 39(4):1127--1139, 1984.
\bibitem{saguillo2025} O.~Saguillo, V.~Ghafouri, L.~Kiffer and G.~Suarez-Tangil. Unravelling the Probabilistic Forest: Arbitrage in Prediction Markets. In \emph{Advances in Financial Technologies (AFT)}, LIPIcs, 2025. arXiv:2508.03474.
\bibitem{snowberg2010} E.~Snowberg and J.~Wolfers. Explaining the favorite-longshot bias: Is it risk-love or misperceptions? \emph{Journal of Political Economy}, 118(4):723--746, 2010.
\bibitem{velickovic2018} P.~Veli\v{c}kovi\'c, G.~Cucurull, A.~Casanova, A.~Romero, P.~Li\`o and Y.~Bengio. Graph attention networks. In \emph{ICLR}, 2018.
\bibitem{wolfers2004} J.~Wolfers and E.~Zitzewitz. Prediction markets. \emph{Journal of Economic Perspectives}, 18(2):107--126, 2004.
\bibitem{yu2018stgcn} B.~Yu, H.~Yin and Z.~Zhu. Spatio-temporal graph convolutional networks: A deep learning framework for traffic forecasting. In \emph{IJCAI}, 2018.
\bibitem{zhang2019deeplob} Z.~Zhang, S.~Zohren and S.~Roberts. DeepLOB: Deep convolutional neural networks for limit order books. \emph{IEEE Transactions on Signal Processing}, 67(11):3001--3012, 2019.
\end{thebibliography}

\end{document}
